\documentclass[10pt,a4paper]{amsart}
\usepackage[margin=19mm]{geometry}
\usepackage{amsmath,amssymb,mathtools}
\usepackage[scr=rsfs,cal=euler]{mathalfa}
\usepackage{xfrac}
\usepackage[unicode,hidelinks,bookmarks=false]{hyperref}
\newcommand{\ie}{i.e.\ }
\newcommand{\cf}{cf.\ }
\newcommand{\resp}{respectively\ }
\newcommand{\Subsection}{\subsection}
\providecommand{\nobreakdash}{\nobreak\hbox{-}}
\setlength{\emergencystretch}{2em}
\multlinegap0pt
\usepackage{xcolor}
\usepackage{amssymb}
\usepackage{algorithm}\usepackage{algpseudocode}
\newtheorem{theorem}{Theorem}
\newcommand{\bA}{\mathbf{A}}
\newcommand{\bB}{\mathbf{B}}
\newcommand{\bF}{\mathbf{F}}
\newcommand{\bM}{\mathbf{M}}
\newcommand{\bo}{\mathbf{0}}
\newcommand{\bu}{\mathbf{u}}
\newcommand{\bw}{\mathbf{w}}
\newcommand{\bx}{\mathbf{x}}
\newcommand{\by}{\mathbf{y}}
\newcommand{\bz}{\mathbf{z}}
\newcommand{\tu}{\tilde{u}}
\newcommand{\tw}{\tilde{w}}
\newcommand{\tz}{\tilde{z}}
\title{A dichotomy of finite element spaces and its application to an energy-conservative scheme for the regularized long wave equation}
\author{Dimitrios Antonopoulos, Dimitrios Mitsotakis}
\date{Source: arXiv:2512.20737v2 (CC BY 4.0)}
\begin{document}
\maketitle

\noindent\emph{Source and licence.} A dichotomy of finite element spaces and its application to an energy-conservative scheme for the regularized long wave equation. Version arXiv:2512.20737v2 (CC BY 4.0), distributed by arXiv under the Creative Commons Attribution 4.0 International licence, http://creativecommons.org/licenses/by/4.0/, as recorded in the arXiv licence metadata for this exact version and shown on the abstract page https://arxiv.org/abs/2512.20737v2. Full licence notice: \texttt{licenses/arxiv-2512.20737-CC-BY-4.0.txt}.\par\smallskip

\noindent\textit{Editorial scope.} Counted unit: the article's theorem thrm:Theorem1 (source0.tex lines 955--958). The article's complete original proof of it is delivered with it (source0.tex lines 960--1102, one \texttt{proof} environment of 143 lines). No mathematics is written, repaired or re-derived: the statement and the proof are the article's own lines, in the article's own notation, and no retained line is substituted or re-flowed.  Retained uncounted as the source-local context the proof needs: lines 603--615; lines 790--792; lines 860--866; lines 890--892; lines 921--942; lines 930--934.  Nothing was omitted from any retained span, so the package carries no omission record.  No retained line contains a citation, so the wrapper carries no bibliography and no bibliography entry.  No retained line contains a figure environment, an image-inclusion command or drawing code, so no image is delivered and the package declares no asset dependency.  Editorial formatting only, in addition: an AMS \texttt{amsart} preamble that restates the article's own declarations and macros used by the retained lines, namely the environments \texttt{theorem} on separate counters, together with the standard packages those lines need.  The retained environments print in this document as Theorem 1 in order of appearance, whereas the article prints the same items with the numbering of its own full document.  The complete native source of the article as distributed in the arXiv e-print for this exact version is bundled unchanged as \texttt{sources/191592/source0.tex}.
\begin{theorem}\label{thm:main1}
\begin{enumerate}
\item[(i)] Let $k=2\kappa+1$, $\kappa=0,1,2,\dots$, and $u\in C_p^{k+2}$. Then there exists a constant $C_k>0$ independent of $h$ and $u$ such that
\begin{equation}\label{eq:onepone}
\|P[(Pu-u)_x]\|\leq C_k h^{k+1}\|u^{(k+2)}\|_\infty\ .
\end{equation}
\item[(ii)]
If $k=1$ and $u\in C_p^5$, then there exists a constant $C_1>0$, independent of $h$ and $u$, such that
\begin{equation}\label{eq:oneptwo}
\|P[(Pu-u)_x]\|\leq C_1 h^4\|u^{(5)}\|_\infty\ .
\end{equation}
\end{enumerate}
\end{theorem}

\begin{equation}\label{eq:inverse}
\|\chi'\|\leq Ch^{-1}\|\chi\|\ .
\end{equation}

\begin{equation}\label{eq:ipbvp}
\begin{aligned}
&u_t+u_x+uu_x-u_{xxt}= 0 \quad\text{for $x\in (a,b)$ and $t>0$}\ ,\\
& \partial^i_x u(a,t)=\partial^i_x u(b,t) \quad \text{ for $t\geq 0$ and $i=0,1,\dots$}\ ,\\
& u(x,0)=U_0(x)\quad \text{for $x\in[a,b]$}\ .
\end{aligned}
\end{equation}

\begin{equation}\label{eq:semid}
(\tu_t,\chi)+(P[\tu_{xt}],\chi_x)=(P[\tu+\tfrac{1}{2}\tu^2],\chi_x)\quad \text{for all $\chi\in \mathcal{V}_h^k$}\ ,
\end{equation}

The semi-discretization (\ref{eq:semid}) can also be expressed in the form:
\begin{alignat}{2}
(\tu_t,\chi)+(\tw_t,\chi_x)&=(\tz,\chi_x) & & \quad \text{for all $\chi\in \mathcal{V}_h^k$}\ ,\label{eq:mixfem1}\\
(\tu_{xt},\psi)-(\tw_t,\psi) &= 0 & & \quad \text{for all $\psi\in \mathcal{V}_h^k$}\ ,\label{eq:mixfem2}\\
(\tz,\phi)&=(\tu+\tfrac{1}{2}\tu^2,\phi)& & \quad \text{for all $\phi\in \mathcal{V}_h^k$}\ .\label{eq:mixfem3}
\end{alignat}
%
In analogy, we can write the RLW equation as
%
\begin{align}
& u_t-w_{xt}=-z_x,\label{eq:bbmn1}\\
& w = u_x, \label{eq:bbmn2}\\
& z = u+\tfrac{1}{2}u^2\ . \label{eq:bbmn3}
\end{align}
%
Note that if $u(x,0)=U_0(x)$, then naturally $w(x,0)=W_0(x)=U_0'(x)$. 
%
We therefore consider as initial conditions of the semi-discrete system (\ref{eq:mixfem1})--(\ref{eq:mixfem3}) the projections
\begin{equation}\label{eq:ics}
\tilde{u}(x,0)=PU_0(x)\quad \text{and}\quad  \tilde{w}(x,0)=PW_0(x)=P[U_0'](x)\ ,
\end{equation}
for all $x\in[0,1]$.

\begin{align}
& u_t-w_{xt}=-z_x,\\
& w = u_x, \\
& z = u+\tfrac{1}{2}u^2\ . 
\end{align}

\begin{theorem}\label{thrm:Theorem1}
Let $u\in C([0,T];C^{k+2}_p)$ be the unique solution of the initial-periodic boundary value problem (\ref{eq:ipbvp}), and let $k=2\kappa+1$ for $\kappa=0,1,\dots$. Then, for sufficiently small $h>0$, there exists a unique solution $\tilde{u}(\cdot,t)\in \mathcal{V}_h^k$ of the semi-discrete problem (\ref{eq:semid}) (and equivalently of (\ref{eq:mixfem1})--(\ref{eq:mixfem3}))  for all $t\in[0,T]$, and a constant $C>0$ independent of $h$ such that
$$\|u-\tilde{u}\|\leq C h^{k+1}, \quad \|u_x-P[\tilde{u}_x]\|\leq C h^{k+1}\quad \text{and}\quad \|u_x-\tilde{u}_x\|\leq C h^k \ .$$
\end{theorem}

\begin{proof}
We define
$$\zeta=u-Pu,\quad \rho=Pu-\tilde{u},\quad \sigma=w-Pw,\quad \theta=Pw-\tilde{w}, \quad\tau=z-Pz,  \quad \xi=Pz-\tilde{z}\ .$$
Note that $\tilde{w}=P[\tilde{u}_x]$ and $\tilde{z}=P[\tilde{u}+\frac{1}{2}\tilde{u}^2]$. 
%
Throughout the analysis, we will use a generic positive constant $C$ independent of $h$.
%
We also assume that the exact solution is bounded, i.e., there is a constant $M$ such that $\|u\|+\|u_x\|\leq M$ for all $t\in[0,T]$. Then, for sufficiently small $h$, the semi-discrete solution satisfies $\|\tilde{u}(\cdot,0)\|+\|\tilde{w}(\cdot,0)\|\leq 2M$.

Upon choosing appropriate basis functions $\varphi_i(x)$, $i=1,2,\dots,kN$ for the finite element space $\mathcal{V}_h^k$, the semi-discrete solution can be represented as a linear combination of the basis functions:
$$\tilde{u}(x,t)=\sum_{i=1}^{kN} u_i(t) \varphi_i(x) \ \text{ and } \ \tilde{w}(x,t)=\sum_{i=1}^{kN} w_i(t)\varphi_i(x)\ .$$
Denote $\by$ the vector $\by=(\bu,\bw)^T$ where $\bu=(u_0,u_1,\dots)^T$ and $\bw=(w_0,w_1,\dots)^T$.
Then, the finite element system (\ref{eq:mixfem1})--(\ref{eq:mixfem3}) can be written as a system of nonlinear ordinary differential equations
\begin{equation}\label{eq:sysnlodes}
   \bM \by' = \bF(\by)\ ,
\end{equation}
where 
$$\bM = \begin{pmatrix}
    \bA & \bB\\
    \bB & \bA
\end{pmatrix}\ ,$$
with $\bA_{ij}=(\varphi_i,\varphi_j)$, $\bB_{ij}=(\varphi_i,\varphi_j')$, and
$\bF=(\bz,\bo)$ with $\bz_j=(\tz,\varphi_j')$. Here, $\bA$ is symmetric and positive definite, while $\bB$ is skew-symmetric. Moreover, for any vector $\bx=(\bx_1, \bx_2)^T$ we have
$$
 \bx^T\bM \bx = \bx_1^T\bA\bx_1+\bx_1^T\bB\bx_2+\bx_2^T\bB\bx_1+\bx_2^T\bA\bx_1> \bx_1^T\bB\bx_2+\bx_1^T(\bx_2^T\bB)^T= 0\ ,
$$
which shows that $\bM$ is positive definite and therefore invertible. 


For the system $\by'=\bM^{-1}\bF(\by)$, the mapping $\bM^{-1}\bF$ is locally Lipschitz as well as $\bF$. Thus, there is a $t^\ast\leq T$ such that the system (\ref{eq:mixfem1})--(\ref{eq:mixfem3}) has a unique solution $(\tilde{u}, \tilde{w})$ satisfying
$\|\tilde{u}\|+\|\tilde{w}\|\leq \tilde{M}$ for all $t\in [0,t^\ast]$, for some constant $\tilde{M}>0$ independent of $h$.

We define $\omega_w=P[w_{xt}]-w_{xt}$, $\omega_z=P[z_x]-z_x$ and $\omega_u=P[u^2]-u^2$. Then, for $\chi\in\mathcal{V}_h^k$, we have 
$$
\begin{aligned}
    (\rho_t,\chi) &= (P[u_t],\chi)-(\tilde{u}_t,\chi)= (P[w_x]_t-P [z_x],\chi) + (\tilde{w}_t-\tilde{z},\chi_x)\\
 &=(P[w_{xt}]-w_{xt},\chi)+(w_{xt},\chi)+(\tilde{w}_t,\chi_x)-(P[z_x]-z_x,\chi)-(z_x,\chi)-(\tilde{z},\chi_x)\\
 &=(\omega_w,\chi)-(w_t-P[w_t],\chi_x)-(P[w_t]-\tilde{w}_t,\chi_x)-(\omega_z,\chi)+(z-Pz,\chi_x)+(Pz-\tilde{z},\chi_x)\\
 &=(P[\omega_w],\chi)+(P[{(\sigma_t)}_x],\chi)+(P[{(\theta_t)}_x],\chi)-(P[\omega_z],\chi)-(P[\tau_x],\chi)-(P[\xi_x],\chi)\ .
\end{aligned}
$$
Equivalently, in operator form, this gives
\begin{equation}\label{eq:rhot}
\rho_t=P[\omega_w+(\sigma_t)_x+(\theta_t)_x-\omega_z-\tau_x-\xi_x]\ .
\end{equation}

From (\ref{eq:mixfem2}) and (\ref{eq:bbmn2}), we have that for any $\psi\in\mathcal{V}_h^k$
$$
\begin{aligned}
(\theta,\psi)&=(Pw-\tilde{w},\psi) = (w-\tilde{w},\psi)= (u_x-P[\tilde{u}_x],\psi)
= -(u-\tilde{u},\psi_x)\\
&= -(u-Pu,\psi_x)-(Pu-\tilde{u},\psi_x)
=  -(\zeta,\psi_x)-(\rho,\psi_x)=  (P[\zeta_x],\psi)+(P[\rho_x],\psi)\ .
\end{aligned}$$
Hence, we obtain the identity
\begin{equation}\label{eq:thetaP}
\theta=P[\zeta_x+\rho_x]\ .
\end{equation}
%
Using the last formula, we have
%
$$
\|P[\rho_x]\| =\|\theta-P[\zeta_x]\|\leq \|\theta\|+\|P[\zeta_x]\|\ .
$$
According to Theorem \ref{thm:main1}, we have $\|P[\zeta_x]\|\leq Ch^{k+1}$. (If $k=1$ and $u\in C_p^5$, the exponent in the last estimate is $4$. However, this won’t improve the error estimate because, for instance, (\ref{eq:xiineq}) below). 
%
Using this estimate, we obtain
\begin{equation}\label{eq:Prhox}
\|P[\rho_x]\|\leq Ch^{k+1}+\|\theta\|\ .
\end{equation}

Let $\phi\in\mathcal{V}_h^k$. Then,
$$
\begin{aligned}
(\xi,\phi) &= (Pz-\tilde{z},\phi)= (Pu-\tilde{u},\phi)+\tfrac{1}{2}(P[u^2]-\tilde{u}^2,\phi)= (\rho,\phi)+\tfrac{1}{2}(P[u^2]-u^2,\phi)+\tfrac{1}{2}(u^2-\tilde{u}^2,\phi)\\
&= (\rho,\phi)+\tfrac{1}{2}(P[u^2]-u^2,\phi)+\tfrac{1}{2}((u+\tilde{u})(u-\tilde{u}),\phi)= (\rho,\phi)+\tfrac{1}{2}(\omega_u,\phi)+\tfrac{1}{2}((u+\tilde{u})(\zeta+\rho),\phi)\ .
\end{aligned}
$$
Taking $\phi=\xi$ in the previous relation yields
$$\|\xi\|^2\leq \|\rho\|\|\xi\|+C h^{k+1}\|\xi\|+\|u+\tilde{u}\|_\infty (C h^{k+1}+\|\rho\|)\|\xi\|\ ,$$
which can be simplified to
\begin{equation}\label{eq:xiineq}
\|\xi\|\leq C(h^{k+1}+\|\rho\|)\ .
\end{equation}
%
Now, we can proceed with the main estimate. From (\ref{eq:rhot}) and (\ref{eq:thetaP}) we have
$$
\begin{aligned}
(\rho_t,\rho)+(\theta_t,\theta) &= (P[\omega_w+(\sigma_t)_x+(\theta_t)_x-\omega_z-\tau_x-\xi_x],\rho)+
(\theta_t, P[\zeta_x+\rho_x]) \\
&=(\omega_w-\omega_z,\rho)+(P[(\sigma_t-\tau)_x],\rho)-(\theta_t,P[\rho_x])-(\xi_x,\rho)+(\theta_t,P[\zeta_x])+(\theta_t,P[\rho_x]) \\
&=(\omega_w-\omega_z,\rho)+(P[(\sigma_t-\tau)_x],\rho)-(\xi_x,\rho)+(\theta_t,P[\zeta_x]) \\
&=(\omega_w-\omega_z,\rho)+(P[(\sigma_t-\tau)_x],\rho)+(\xi,P[\rho_x])+\tfrac{d}{dt}(\theta,P[\zeta_x])-(\theta,P[(\zeta_t)_x])
\ .
\end{aligned}
$$
The last term in the previous inequality can be bounded using the Cauchy-Schwarz inequality: 
$$
(\theta,P[(\zeta_t)_x])\leq \frac{1}{8}\|\theta\|^2+2\|P[(\zeta_t)_x]\|^2\leq C(\|P[(\zeta_t)_x]\|^2+\|\theta\|^2)\ .
$$
%
By Theorem \ref{thm:main1}, $\|P[(\zeta_t)_x]\|\leq C h^{k+1}$. Therefore,
\begin{equation}\label{eq:estim1}
(\theta,P[(\zeta_t)_x])\leq C( h^{2(k+1)} +\|\theta\|^2) \  .
\end{equation}
%
Similarly, using the inequalities (\ref{eq:xiineq}) and (\ref{eq:Prhox}), we obtain
\begin{equation}\label{eq:estim2}
(\xi,P[\rho_x]) \leq C(h^{2(k+1)}+\|\xi\|^2+\|\theta\|^2)\leq C(h^{2(k+1)} +\|\rho\|^2+\|\theta\|^2)\ .
\end{equation}
%
Thus,
\begin{equation}
\frac{1}{2}\frac{d}{dt}(\|\rho\|^2+\|\theta\|^2) \leq C(h^{2(k+1)}+\|\rho\|^2+\|\theta\|^2)+\tfrac{d}{dt}(\theta,P[\zeta_x]) \ .   
\end{equation}
%
Integrating this relation yields 
\begin{equation}\label{eq:lrel}
\|\rho\|^2+\|\theta\|^2\leq C h^{2(k+1)} + \frac{1}{8}\|\theta\|^2+2\|P[\zeta_x]\|^2 + C\int_0^t (\|\rho\|^2+\|\theta^2\|)\ ,
\end{equation}
where, from now on, $C$ depends on $t^\ast$.
Taking into account that $\|P[\zeta_x]\|\leq C h^{k+1}$, we can simplify (\ref{eq:lrel}) to
\begin{equation}
\|\rho\|^2+\|\theta\|^2\leq Ch^{2(k+1)}+C\int_0^t (\|\rho\|^2+\|\theta^2\|)\ .
\end{equation}
%
Applying Gronwall's inequality, we obtain the optimal error estimate
\begin{equation}\label{eq:help2}
\|\rho\|^2+\|\theta\|^2\leq Ch^{2(k+1)}\ ,
\end{equation}
for all $t\in [0,t^\ast]$, which immediately gives
\begin{equation}
\|u-\tilde{u}\|^2+\|u_x-P[\tilde{u}_x]\|^2\leq Ch^{2(k+1)}\ .
\end{equation}
%
Finally, using (\ref{eq:help2}) and the inverse inequality (\ref{eq:inverse}), we have 
$$\|u_x-\tilde{u}_x\|\leq \|(u-Pu)_x\|+\|\rho_x\|\leq C(h^k+h^{-1}\|\rho\|)\leq Ch^k\ .$$
%
Choosing $h>0$ sufficiently small such that $\|\tilde{u}\|\leq \|u\|+Ch^{k+1}<2M$ for $0\leq t\leq t^\ast$, 
we ensure that the numerical solution $\tilde{u}$ will remain bounded in $L^2$ up to $t=t^\ast$. Similarly, since $\|u_x-P[\tilde{u}_x]\|\leq Ch^{k+1}$, we have 
$\|\tilde{w}\|\leq \|u_x\|+Ch^{k+1}<2M$ for $0\leq t\leq t^\ast$.
By standard extension arguments from the theory of ordinary differential equations, we conclude that the solution (and the error estimates) can be extended to the full interval $t\in[0,T]$. 
\end{proof}

\end{document}
